% Extracto íntegro por residencia del cierre físico-matemático sucesor.
\section{Entrelazador canónico entre memoria TPK y aparato}
\label{sec:hmtf2-aparato-memoria}

Sea \(\mathfrak M_n\) el conjunto finito de registros TPK a nivel \(n\),
incluidos resultado, ruta, acarreo, hoja y memoria. Su Hilbertización
canónica es
\[
 \mathcal H_{\mathfrak M_n}
 =\ell^2(\mathfrak M_n).
\]
Sea \(\mathcal K_{A,n}\) un espacio auxiliar del aparato que contiene una
familia ortonormal
\(\{|a_m\rangle:m\in\mathfrak M_n\}\). Definimos
\[
 \iota_n:\mathcal H_{\mathfrak M_n}\longrightarrow\mathcal K_{A,n},
 \qquad
 \iota_n|m\rangle=|a_m\rangle.
\]
Para no perder genealogía cuando la actualización visible no sea inyectiva,
se usa el levantamiento de ledger
\[
 \widehat F_n(m)=(F_n(m),m)\in\mathfrak M_{n+1}.
\]
Es inyectivo por construcción y define la isometría
\[
 V_n|m\rangle=|\widehat F_n(m)\rangle.
\]
Sobre el subespacio de puntero se fija
\[
 W_n\iota_n=\iota_{n+1}V_n.
\]

\begin{theorem}[Entrelazador aparato--memoria TPK]
\label{thm:hmtf2-entrelazador-memoria}
Los mapas \(\iota_n\) y \(V_n\) son isometrías y el cuadrado
\[
 \begin{array}{ccc}
 \mathcal H_{\mathfrak M_n}&\xrightarrow{\ V_n\ }&
 \mathcal H_{\mathfrak M_{n+1}}\\
 {\scriptstyle\iota_n}\downarrow&&\downarrow{\scriptstyle\iota_{n+1}}\\
 \mathcal K_{A,n}&\xrightarrow{\ W_n\ }&
 \mathcal K_{A,n+1}
 \end{array}
\]
conmuta. El mapa \(W_n\) se prolonga a una isometría del aparato y, tras
añadir un complemento auxiliar si las dimensiones difieren, a una unitaria.
La actualización preserva norma, distingue todos los registros y conserva el
estado anterior dentro del ledger.
\end{theorem}

\begin{proof}
Los mapas envían bases ortonormales en familias ortonormales. La identidad
del cuadrado se verifica sobre cada \(|m\rangle\). Toda isometría finita se
extiende a una unitaria entre dilataciones de igual dimensión completando
bases ortonormales.
\end{proof}

Los proyectores de puntero
\[
 P_m=|a_m\rangle\langle a_m|
\]
definen la expectativa condicional
\[
 \mathcal D_n(A)=\sum_{m\in\mathfrak M_n}P_m A P_m.
\]
Se cumple
\[
 \mathcal D_n^2=\mathcal D_n,\qquad
 \operatorname{Spec}(\mathcal D_n)\subset\{0,1\}.
\]
La subálgebra diagonal de registros es el subespacio fijo y el complemento
coherente pertenece al núcleo. El gap unitario de esta proyección vale uno.
Bajo una conjugación simultánea del aparato, de los punteros y del
entrelazador, todas estas identidades permanecen exactas. Así queda cerrada
la estabilidad finita del registro dentro de la clase covariante.

Cuando los rótulos \((y,\alpha)\) de Kraus se incluyen en
\(\mathfrak M_n\), la isometría \(\iota_n\) identifica explícitamente la
memoria TPK con el ambiente de Stinespring, no sólo con un espacio auxiliar
abstracto. La construcción de un dispositivo material, su decoherencia y
su calibración constituyen \textsc{reconocimiento externo}.


